\FigureLayoutDeclare{img/2020/national_paper_3_science/q8_solution_fig1.png}{4.0cm}{6.999bp 0bp 28.996bp 1bp}
\FigureLayoutDeclare{img/2020/national_paper_3_science/q15_solution_fig1.png}{5.8cm}{0bp 0bp 5.999bp 0bp}

\chapter{2020年全国III卷（理）}

\section{单选题}

\begin{problem}
已知集合 \(A=\{(x,y)\mid x,y\in \mathbf N^*,y\ge x\}\)，\(B=\{(x,y)\mid x+y=8\}\)，则 \(A\cap B\) 中元素的个数为（\quad）

\choices
    {\(2\)}
    {\(3\)}
    {\(4\)}
    {\(6\)}
\end{problem}

\begin{problem}
复数 \(\dfrac{1}{1-3\i}\) 的虚部是（\quad）

\choices
    {\(-\dfrac{3}{10}\)}
    {\(-\dfrac{1}{10}\)}
    {\(\dfrac{1}{10}\)}
    {\(\dfrac{3}{10}\)}
\end{problem}

\begin{problem}
在一组样本数据中，\(1\)，\(2\)，\(3\)，\(4\) 出现的频率分别为 \(p_1\)，\(p_2\)，\(p_3\)，\(p_4\)，且 \(\sum_{i=1}^{4} p_i=1\)，则下面四种情形中，对应样本的标准差最大的一组是（\quad）

\choices
    {\(p_1=p_4=0.1\)，\(p_2=p_3=0.4\)}
    {\(p_1=p_4=0.4\)，\(p_2=p_3=0.1\)}
    {\(p_1=p_4=0.2\)，\(p_2=p_3=0.3\)}
    {\(p_1=p_4=0.3\)，\(p_2=p_3=0.2\)}
\end{problem}

\begin{problem}
Logistic 模型是常用数学模型之一，可应用于流行病学领域. 有学者根据公布数据建立了某地区新冠肺炎累计确诊病例数 \(I(t)\)（\(t\) 的单位：天）的 Logistic 模型：\(I(t)=\dfrac{K}{1+\e^{-0.23(t-53)}}\)，其中 \(K\) 为最大确诊病例数. 当 \(I(t^*)=0.95K\) 时，标志着已初步遏制疫情，则 \(t^*\) 约为（\quad）（\(\ln19\approx3\)）

\choices
    {\(60\)}
    {\(63\)}
    {\(66\)}
    {\(69\)}
\end{problem}

\begin{problem}
设 \(O\) 为坐标原点，直线 \(x=2\) 与抛物线 \(C:y^2=2px\)（\(p>0\)）交于 \(D\)，\(E\) 两点，若 \(OD\perp OE\)，则 \(C\) 的焦点坐标为（\quad）

\choices
    {\(\left( \dfrac{1}{4},0 \right)\)}
    {\(\left( \dfrac{1}{2},0 \right)\)}
    {\((1,0)\)}
    {\((2,0)\)}
\end{problem}

\begin{problem}
已知向量 \(\bs{a}\)，\(\bs{b}\) 满足 \(|\bs{a}|=5\)，\(|\bs{b}|=6\)，\(\bs{a}\cdot\bs{b}=-6\)，则 \(\cos\angle(\bs{a},\bs{a}+\bs{b})=\)（\quad）

\choices
    {\(-\dfrac{31}{35}\)}
    {\(-\dfrac{19}{35}\)}
    {\(\dfrac{17}{35}\)}
    {\(\dfrac{19}{35}\)}
\end{problem}

\begin{problem}
在 \(\triangle ABC\) 中，\(\cos C=\dfrac{2}{3}\)，\(AC=4\)，\(BC=3\)，则 \(\cos B=\)（\quad）

\choices
    {\(\dfrac{1}{9}\)}
    {\(\dfrac{1}{3}\)}
    {\(\dfrac{1}{2}\)}
    {\(\dfrac{2}{3}\)}
\end{problem}

\begin{problem}
如图为某几何体的三视图，则该几何体的表面积是（\quad）

\bitmapfigure[max width=6.0cm]{img/2020/national_paper_3_science/q8_fig1.png}

\choices
    {\(6+4\sqrt{2}\)}
    {\(4+4\sqrt{2}\)}
    {\(6+2\sqrt{3}\)}
    {\(4+2\sqrt{3}\)}
\end{problem}

\begin{problem}
已知 \(2\tan\theta-\tan\left( \theta+\dfrac{\pi}{4} \right)=7\)，则 \(\tan\theta=\)（\quad）

\choices
    {\(-2\)}
    {\(-1\)}
    {\(1\)}
    {\(2\)}
\end{problem}

\begin{problem}
若直线 \(l\) 与曲线 \(y=\sqrt{x}\) 和圆 \(x^2+y^2=\dfrac{1}{5}\) 都相切，则 \(l\) 的方程为（\quad）

\choices
    {\(y=2x+1\)}
    {\(y=2x+\dfrac{1}{2}\)}
    {\(y=\dfrac{1}{2}x+1\)}
    {\(y=\dfrac{1}{2}x+\dfrac{1}{2}\)}
\end{problem}

\begin{problem}
设双曲线 \(C:\dfrac{x^2}{a^2}-\dfrac{y^2}{b^2}=1\)（\(a>0\)，\(b>0\)）的左，右焦点分别为 \(F_1\)，\(F_2\)，离心率为 \(\sqrt{5}\). \(P\) 是 \(C\) 上一点，且 \(F_1P\perp F_2P\). 若 \(\triangle PF_1F_2\) 的面积为 \(4\)，则 \(a=\)（\quad）

\choices
    {\(1\)}
    {\(2\)}
    {\(4\)}
    {\(8\)}
\end{problem}

\begin{problem}
已知 \(5^5<8^4\)，\(13^4<8^5\). 设 \(a=\log_5 3\)，\(b=\log_8 5\)，\(c=\log_{13} 8\)，则（\quad）

\choices
    {\(a<b<c\)}
    {\(b<a<c\)}
    {\(b<c<a\)}
    {\(c<a<b\)}
\end{problem}

\section{填空题}

\begin{problem}
若 \(x\)，\(y\) 满足 \(\begin{cases}
x+y\ge0\\
2x-y\ge0\\
x\le1
\end{cases}\)
，则 \(z=3x+2y\) 的最大值为\fillinblank{}.
\end{problem}

\begin{problem}
\(\left( x^2+\dfrac{2}{x} \right)^6\) 的展开式中常数项是\fillinblank{}（用数字作答）.
\end{problem}

\begin{problem}
已知圆锥的底面半径为 \(1\)，母线长为 \(3\)，则该圆锥内半径最大的球的体积为\fillinblank{} \(\pi\).
\end{problem}

\begin{problem}
关于函数 \(f(x)=\sin x+\dfrac{1}{\sin x}\) 有如下四个命题：

\begin{circlelist}
\item \(f(x)\) 的图象关于 \(y\) 轴对称；
\item \(f(x)\) 的图象关于原点对称；
\item \(f(x)\) 的图象关于直线 \(x=\dfrac{\pi}{2}\) 对称；
\item \(f(x)\) 的最小值为 \(2\).
\end{circlelist}

其中所有真命题的序号是\fillinblank{}.
\end{problem}

\section{解答题}

\begin{problem}
设数列 \(\{a_n\}\) 满足 \(a_1=3\)，\(a_{n+1}=3a_n-4n\).

\begin{enumerate}
\item 计算 \(a_2\)，\(a_3\)，猜想 \(\{a_n\}\) 的通项公式并加以证明；
\item 求数列 \(\{2^na_n\}\) 的前 \(n\) 项和 \(S_n\).
\end{enumerate}
\end{problem}

\begin{problem}
某学生兴趣小组随机调查了某市 \(100\) 天中每天的空气质量等级和当天到某公园锻炼的人次，整理数据得到下表（单位：天）：

\begin{center}
\begin{tabular}{c|ccc}
锻炼人次 & \([0,200]\) & \((200,400]\) & \((400,600]\) \\
\hline
空气质量等级 \(1\)（优） & 2 & 16 & 25 \\
空气质量等级 \(2\)（良） & 5 & 10 & 12 \\
空气质量等级 \(3\)（轻度污染） & 6 & 7 & 8 \\
空气质量等级 \(4\)（中度污染） & 7 & 2 & 0
\end{tabular}
\end{center}

\begin{enumerate}
\item 分别估计该市一天的空气质量等级为 \(1\)，\(2\)，\(3\)，\(4\) 的概率；
\item 求一天中到该公园锻炼的平均人次的估计值（同一组中的数据用该组区间的中点值为代表）；
\item 若某天的空气质量等级为 \(1\) 或 \(2\)，则称这天“空气质量好”；若某天的空气质量等级为 \(3\) 或 \(4\)，则称这天“空气质量不好”. 根据所给数据，完成下面的 \(2\times2\) 列联表，并根据列联表，判断是否有 \(95\%\) 的把握认为一天中到该公园锻炼的人次与该市当天的空气质量有关？
\end{enumerate}

\begin{center}
\begin{tabular}{c|cc}
 & 人次\(\le400\) & 人次\(>400\) \\
\hline
空气质量好 &  &  \\
空气质量不好 &  &
\end{tabular}
\end{center}

附：\(K^2=\dfrac{n(ad-bc)^2}{(a+b)(c+d)(a+c)(b+d)}\).

\begin{center}
\begin{tabular}{c|ccc}
\(P(K^2\ge k)\) & 0.050 & 0.010 & 0.001 \\
\hline
\(k\) & 3.841 & 6.635 & 10.828
\end{tabular}
\end{center}
\end{problem}

\begin{problem}
如图，在长方体 \(ABCD-A_1B_1C_1D_1\) 中，点 \(E\)，\(F\) 分别在棱 \(DD_1\)，\(BB_1\) 上，且 \(2DE=ED_1\)，\(BF=2FB_1\).

\begin{enumerate}
\item 证明：点 \(C_1\) 在平面 \(AEF\) 内；
\item 若 \(AB=2\)，\(AD=1\)，\(AA_1=3\)，求二面角 \(A-EF-A_1\) 的正弦值.
\end{enumerate}

\bitmapfigure[max width=6.0cm]{img/2020/national_paper_3_science/q19_fig1.png}
\end{problem}

\begin{problem}
已知椭圆 \(C:\dfrac{x^2}{25}+\dfrac{y^2}{m^2}=1\)（\(0<m<5\)）的离心率为 \(\dfrac{\sqrt{15}}{4}\)，\(A\)，\(B\) 分别为 \(C\) 的左，右顶点.

\begin{enumerate}
\item 求 \(C\) 的方程；
\item 若点 \(P\) 在 \(C\) 上，点 \(Q\) 在直线 \(x=6\) 上，且 \(|BP|=|BQ|\)，\(BP\perp BQ\)，求 \(\triangle APQ\) 的面积.
\end{enumerate}
\end{problem}

\begin{problem}
设函数 \(f(x)=x^3+bx+c\)，曲线 \(y=f(x)\) 在点 \(\left( \dfrac{1}{2},f\left( \dfrac{1}{2} \right) \right)\) 处的切线与 \(y\) 轴垂直.

\begin{enumerate}
\item 求 \(b\)；
\item 若 \(f(x)\) 有一个绝对值不大于 \(1\) 的零点，证明：\(f(x)\) 所有零点的绝对值都不大于 \(1\).
\end{enumerate}
\end{problem}

\section{选考题：解答题}

\begin{problem}
在直角坐标系 \(xOy\) 中，曲线 \(C\) 的参数方程为 \(x=2-t-t^2\)，\(y=2-3t+t^2\)（\(t\) 为参数且 \(t\ne1\)），\(C\) 与坐标轴交于 \(A\)，\(B\) 两点.

\begin{enumerate}
\item 求 \(|AB|\)；
\item 以坐标原点为极点，\(x\) 轴正半轴为极轴建立极坐标系，求直线 \(AB\) 的极坐标方程.
\end{enumerate}
\end{problem}

\begin{problem}
设 \(a\)，\(b\)，\(c\in \R\)，\(a+b+c=0\)，\(abc=1\).

\begin{enumerate}
\item 证明：\(ab+bc+ca<0\)；
\item 用 \(\max\{a,b,c\}\) 表示 \(a\)，\(b\)，\(c\) 中的最大值，证明：\(\max\{a,b,c\}\ge \sqrt[3]{4}\).
\end{enumerate}
\end{problem}

